\documentclass[11pt]{article}

\usepackage[margin=1.3in]{geometry}
\usepackage[utf8]{inputenc}
\usepackage[UKenglish]{babel}

\usepackage{amsmath}
\usepackage{amsthm}
\usepackage{amsfonts}
\usepackage{amssymb}
\usepackage{hyperref}

\newtheorem{theorem}{Theorem}
\newtheorem{lemma}{Lemma}
\newtheorem{corollary}{Corollary}
\newtheorem{definition}{Definition}
\newtheorem{example}{Example}

\newcommand{\eps}{\varepsilon}
\newcommand{\R}{\mathbb{R}}
\newcommand{\Z}{\mathbb{Z}}
\newcommand{\N}{\mathcal{N}}
\newcommand{\E}{\mathrm{E}}
\newcommand{\tr}{\mathrm{trace}}

\begin{document}

\section{Simple Stirling Upper/Lower Bounds}

Stirling's approximation is often proved with heavy machinery, but there's a short, elementary route starting from the integral $n! = \int_0^\infty x^n e^{-x} \, dx$. With a single change of variables that turns $\log(1+x) - x$ into $-z^{2}/2$, the problem becomes Gaussian integration: the main mass contributes $\sqrt{2\pi n}$, and the remainder is controlled by simple logarithmic inequalities. In a few lines this yields the sharp bounds $\left(\frac{n}{e}\right)^n\sqrt{2\pi n} \le n! \le \left(\frac{n}{e}\right)^n\bigl(\sqrt{2\pi n}+1\bigr)$. The same viewpoint also exposes the higher-order corrections by integrating low-degree polynomials against the Gaussian, recovering the familiar $1/(12n)$, $1/(288n^2)$, $\ldots$ terms. Originally posted here: \href{https://mathoverflow.net/a/484270/5429}{mathoverflow.net/a/484270/5429}.

We will give a simple proof for the bound
\[
 (n/e)^n\sqrt{2\pi n}
 \le
 n! \le (n/e)^n(\sqrt{2\pi n}+1).
\]

We use the integral representation of $n!$:
\begin{align*}
    n!
    &= \int_0^\infty x^n e^{-x} \\
    &= n\int_0^\infty e^{n \log(nx) - nx} \, dx \\
    &= n n^{n}\int_0^\infty e^{n (\log(x) - x)} \, dx \\
    &= n (n/e)^{n}\int_{-1}^\infty e^{n (\log(1+x) - x)} \, dx.
\end{align*}
Note that $\log(1+x)-x$ is maximized at $x=0$.
Inspired by the series expansion $\log(1+x)-x=-x^2/2 + O(x^3)$ we make a change of variables,
$\log(1+x)-x=-z^2/2$, or:
\[
z = \mathrm{sign}(x)\,\sqrt{2\bigl(x-\log(1+x)\bigr)}.
\]
The differentials are simple:
\[
d(-z^2/2)=
-z\,dz = d(\log(1+x)-x) = \frac{dx}{1+x} - dx = -\frac{x}{1+x}dx,
\]
which after adjusting the integration limits yield:
\[
    n!
    = (n/e)^{n} \, n\int_{-\infty}^\infty e^{-nz^2/2} \, \bigl(z+z/x\bigr) \, dz.
\]
We now need the bounds
\begin{equation}
    1+\tfrac{2}{3}z
    \;\le\; z+\frac{z}{x} \;\le\; \max\{1, 1+z\}
    \quad\text{for all }x\in\R.
    \label{eq:mid}
\end{equation}
Since $\int_{-\infty}^\infty e^{-nz^2/2}dz=\sqrt{2\pi/n}$ and $\int_{0}^\infty e^{-nz^2/2} z\, dz=1/n$,
this immediately gives us
\[
\sqrt{2\pi n}
\le
n\int_{-\infty}^\infty e^{-nz^2/2} \, (z+z/x) \, dz
\le
\sqrt{2\pi n} + 1.
\]
Which is what we want.

It remains to show \eqref{eq:mid},
which we will do using the following standard logarithmic inequalities~\cite{wiki-logs}:
\begin{align}
    \log(1+x) &\ge \frac{x(2+x)}{2+2x} \quad &&\text{when } x<0
    \label{eq:logneg}
    \\
    \log(1+x) &\ge x-\tfrac{1}{2}x^2 \quad &&\text{when } x>0
    \label{eq:loglower}
    \\
    \log(1+x) &\le \frac{x(6+x)}{6+4x}
     \quad &&\text{when } x>-1.
    \label{eq:logupper}
\end{align}
We note $z(1+1/x)$ is non-negative, since $1+1/x<0$ implies $\mathrm{sign}(x)=-1$.
Thus it suffices to show $z^2(1+1/x)^2\le 1$, which follows using \eqref{eq:logneg}:
\[
z^2 = 2\bigl(x-\log(1+x)\bigr)
\le \frac{x^2}{1+x}
\le \frac{x^2}{(1+x)^2}
= (1+1/x)^{-2},
\]
For $x>0$ we can use \eqref{eq:loglower} to bound
\[
z = \sqrt{2\bigl(x-\log(1+x)\bigr)}
\le \sqrt{2(x^2/2)} = x,
\]
so $z/x\le 1$.
This proves the upper bound of \eqref{eq:mid}.
For the lower bound, it suffices to show
$z^2(1/3+1/x)^2 \ge 1.$
Using \eqref{eq:logupper} we get
\begin{align*}
z^2\Bigl(\tfrac{1}{3}+\tfrac{1}{x}\Bigr)^2
&= 2\bigl(x - \log(1+x)\bigr)\Bigl(\tfrac{1}{3}+\tfrac{1}{x}\Bigr)^2 \\
&\ge 2\Bigl(x - \frac{x(6+x)}{6+4x}\Bigr)\Bigl(\tfrac{1}{3}+\tfrac{1}{x}\Bigr)^2 \\
&= \frac{x^2}{9 + 6 x} + 1
\quad\ge 1 \quad \text{since } x\ge -1.
\end{align*}

\medskip

The full Stirling approximation can similarly be derived from the series expansion:
\[
z\Bigl(1+\frac{1}{x}\Bigr) = 1 + \frac{2z}{3} + \frac{z^2}{12}
- \frac{2z^3}{135} + \frac{z^4}{864} + \dots
\]
This form is nice as it allows easy integration with respect to the Gaussian pdf.
E.g.
\[
n\int_{-\infty}^{\infty}\frac{z^4}{864} \, e^{-n z^2/2} \, dz = \frac{\sqrt{2 \pi n}}{288n^{2}}
\]
matching the expected
\[
n! \sim \sqrt{2\pi n}\left(\frac{n}{e}\right)^n \left(1 +\frac{1}{12n}+\frac{1}{288n^2} - \frac{139}{51840n^3} -\frac{571}{2488320n^4}+ \cdots \right).
\]

I think one can even give a uniform bound of
\[
z\Bigl(1+\frac{1}{x}\Bigr) \le 1 + \frac{2z}{3} + \frac{z^2}{c},
\]
for some $c\approx 8$, strengthening the upper bound even more.
However, the logarithmic inequalities become quite tedious.

\begin{thebibliography}{1}
\bibitem{wiki-logs}
Wikipedia contributors, ``List of logarithmic identities — Inequalities,''
\emph{Wikipedia}, \url{https://en.wikipedia.org/wiki/List_of_logarithmic_identities#Inequalities}.
\end{thebibliography}

\end{document}
